
\chapter{Hypothesis Spaces and Inductive Bias}

\section{Function Classes and Model Families}

In previous chapters, we formulated learning as the problem of selecting a function that minimizes empirical risk. However, we have not yet addressed a crucial question:

\medskip

\noindent
\textbf{Which functions are we allowed to choose from?}

\medskip

This question is answered by specifying a \emph{hypothesis space} (or model class)
\[
\mathcal{H} = \{ f_\theta : \theta \in \Theta \}.
\]

\medskip

The hypothesis space defines the set of candidate functions that a learning algorithm can represent.

\medskip

\noindent
\textbf{Examples of function classes.}

\begin{itemize}
    \item \textbf{Linear models:}
    \[
    f(x) = \langle w, x \rangle + b.
    \]

    \item \textbf{Polynomial models:}
    \[
    f(x) = \sum_{|\alpha| \leq d} c_\alpha x^\alpha.
    \]

    \item \textbf{Neural networks:}
    \[
    f_\theta(x) = W_L \sigma\big(W_{L-1} \sigma(\cdots \sigma(W_1 x))\big).
    \]
\end{itemize}

\medskip

Each choice of $\mathcal{H}$ imposes structural assumptions about the relationship between inputs and outputs.

\medskip

\noindent
\textbf{Expressivity.}  
A hypothesis space is said to be expressive if it can approximate a wide range of functions.

\medskip

\noindent
\textbf{Example.}  
Neural networks with sufficient width can approximate any continuous function on compact domains (universal approximation property).

\medskip

\noindent
However, expressivity alone does not determine performance. A highly expressive model may fit training data perfectly but fail to generalize.

\medskip

Thus, the choice of $\mathcal{H}$ is not merely about approximation—it is about balancing expressivity, stability, and learnability.

\section{Inductive Bias and Prior Structure}

The hypothesis space introduces an implicit preference for certain functions over others. This preference is known as \emph{inductive bias}.

\medskip

\noindent
\textbf{Definition (informal).}  
Inductive bias is the set of assumptions that allows a learning algorithm to generalize beyond the observed data.

\medskip

\noindent
Without inductive bias, learning is impossible: infinitely many functions can interpolate the training data, but only some will generalize well.

\medskip

\noindent
\textbf{Sources of inductive bias.}

\begin{itemize}
    \item \textbf{Model class:}  
    Linear models favor linear relationships.

    \item \textbf{Architecture:}  
    Convolutional networks favor locality and translation invariance.

    \item \textbf{Regularization:}  
    Penalizing large parameters favors simpler functions.

    \item \textbf{Optimization:}  
    Gradient-based methods may prefer certain solutions among many.
\end{itemize}

\medskip

\noindent
\textbf{Connection to prior knowledge.}  
Inductive bias can be viewed as encoding prior knowledge about the problem:
\begin{itemize}
    \item smoothness,
    \item sparsity,
    \item compositional structure,
    \item invariances.
\end{itemize}

\medskip

\noindent
\textbf{Engineering perspective.}  
Successful models are not neutral—they are carefully designed to reflect structure in the data.

\medskip

\noindent
For example:
\begin{itemize}
    \item CNNs exploit spatial locality in images,
    \item Transformers exploit relationships between elements in sequences.
\end{itemize}

\medskip

Inductive bias is therefore not a limitation, but a necessary ingredient for learning.

\section{Bias--Variance Tradeoff (Intuition)}

We now examine how the choice of hypothesis space affects generalization.

\medskip

\noindent
Consider the expected prediction error at a point $x$:
\[
\mathbb{E}\big[(f(x) - y)^2\big].
\]

\medskip

This error can be conceptually decomposed into three components:

\begin{itemize}
    \item \textbf{Bias:} error due to systematic approximation
    \item \textbf{Variance:} error due to sensitivity to data
    \item \textbf{Noise:} inherent randomness in the data
\end{itemize}

\medskip

\noindent
\textbf{High-bias models.}
\begin{itemize}
    \item Simple hypothesis spaces (e.g., linear models)
    \item May underfit the data
    \item Low variance but large systematic error
\end{itemize}

\medskip

\noindent
\textbf{High-variance models.}
\begin{itemize}
    \item Highly expressive models (e.g., deep networks)
    \item Can fit training data closely
    \item Sensitive to noise and fluctuations
\end{itemize}

\medskip

\noindent
\textbf{Tradeoff.}  
Increasing model complexity typically:
\begin{itemize}
    \item decreases bias,
    \item increases variance.
\end{itemize}

\medskip

\noindent
\textbf{Insight.}  
Good generalization requires balancing these two effects.

\medskip

\noindent
\textbf{Remark.}  
In modern deep learning, this classical picture is refined: highly overparameterized models can still generalize well due to implicit regularization and optimization effects.

\section{Choosing Representations}

The performance of a learning system depends not only on the function class, but also on how inputs are represented.

\medskip

\noindent
\textbf{Definition.}  
A representation is a transformation
\[
\phi : \mathcal{X} \to \mathcal{Z}
\]
that maps raw inputs into a feature space.

\medskip

A model then operates on $\phi(x)$ instead of $x$:
\[
f(x) = g(\phi(x)).
\]

\medskip

\noindent
\textbf{Example.}
\begin{itemize}
    \item Polynomial features: $\phi(x) = (x, x^2, \dots)$
    \item Learned features in neural networks
\end{itemize}

\medskip

\noindent
\textbf{Key idea.}  
A simple model applied to a rich representation can be more powerful than a complex model applied to raw inputs.

\medskip

\noindent
\textbf{Feature engineering vs representation learning.}

\begin{itemize}
    \item Classical approaches rely on manually designed features
    \item Deep learning learns representations automatically
\end{itemize}

\medskip

\noindent
\textbf{Compositional representations.}  
Deep networks build representations hierarchically:
\begin{itemize}
    \item lower layers capture local patterns,
    \item higher layers capture abstract structure.
\end{itemize}

\medskip

\noindent
\textbf{Geometric viewpoint.}  
Representations transform the geometry of the data so that:
\begin{itemize}
    \item classes become separable,
    \item structure becomes simpler,
    \item learning becomes easier.
\end{itemize}

\medskip

Thus, choosing a representation is equivalent to choosing a coordinate system in which the learning problem becomes tractable.

\section{A Unifying Perspective}

We can now refine the learning problem:

\begin{enumerate}
    \item Choose a representation $\phi$
    \item Choose a hypothesis space $\mathcal{H}$
    \item Impose inductive bias through structure and regularization
    \item Optimize empirical risk
\end{enumerate}

\medskip

\noindent
\textbf{Key components of learning:}
\begin{itemize}
    \item \textbf{Representation} (how data is encoded)
    \item \textbf{Function class} (what can be expressed)
    \item \textbf{Inductive bias} (what is preferred)
\end{itemize}

\medskip

Modern deep learning systems unify these elements: they learn representations and functions simultaneously, guided by implicit and explicit biases.

\section{Outlook}

This chapter emphasized that learning is not only about fitting data, but about choosing the right space of functions and representations.

\medskip

In the next chapter, we study learning as an optimization problem, focusing on how models are trained in practice.

\medskip

\noindent
\textbf{Guiding principle.}  
Generalization arises not from data alone, but from the interaction between data, model structure, and inductive bias.
